% vim: ai sw=2 spell spelllang=cs

\begin{frame}{Knihovny}

  \begin{itemize}
    \item Statické knihovny
    \begin{itemize}
      \item \cmd{ar} archiv
      \item Jejích \emph{obsah} je v každé binárce, kam se linkovaly
      \item Zvyšují velikost na disku
      \item Binárka ale nemá žádné externí závislosti

      \pause

      \item Výroba: \cmd{gcc -c} a \cmd{ar rcs lib.a *.o}
    \end{itemize}

    \pause

    \item Dynamické knihovny
    \begin{itemize}
      \item ELF, který se linkuje za chodu
      \begin{itemize}
	\item Podpora pluginů
	\item Ale zpomaluje spuštění

	\pause

	\item Dynamický linker \file{ld.so} (\doc{man ld.so})
      \end{itemize}

      \pause

      \item Výroba: \cmd{gcc \-shared}
    \end{itemize}

  \end{itemize}

\end{frame}

\iffimuni
\begin{frame}{Úkol}

  \heading{Výroba statické knihovny}
  \begin{enumerate}
    \item Proveďte \cmd{make} v \file{pb173-bin/12}

    \item Z \file{lib.o} vyrobte statickou knihovnu \file{libX.a}
    \begin{itemize}
      \item \cmd{ar rcs ...}
    \end{itemize}

    \item Knihovnu přilinkujte k \file{main.o}
    \begin{itemize}
      \item \cmd{gcc -L. -lX}
    \end{itemize}

    \item Spusťte výsledek
  \end{enumerate}

\end{frame}
\fi

\begin{frame}{Použití dynamické knihovny}

  \begin{enumerate}
    \item Přilinkování pomocí linkeru
    \begin{itemize}
      \item \cmd{gcc -lknihovna}
      \item \file{ld.so} před spuštěním načte a slinkuje potřebné symboly

      \pause

      \item \heading{Demo:} \cmd{ldd} k ověření
    \end{itemize}

    \pause

    \item Načtení explicitně až za běhu
    \begin{itemize}
      \item Viz dále
    \end{itemize}

    \pause

    \item Proměnnými prostředí
    \begin{itemize}
      \item \cmd{LD_PRELOAD} nahraje knihovnu před spuštěním
    \end{itemize}
  \end{enumerate}

\end{frame}

\begin{frame}{Načtení dynamické knihovny}

  \begin{itemize}
    \item Knihovna \file{dl} (linkujte s \cmd{\-ldl}) (\header{dlfcn.h})

    \pause

    \item Otevření: \code{void *dlopen(const char *filename, int flags)}
    \begin{itemize}
      \item \code{filename}: jméno knihovny, nebo cesta
      \item \code{flags}: \code{RTLD_LAZY} nebo \code{RTLD_NOW} spolu s
        dalšími
      \item Návratová hodnota: držátko knihovny nebo \code{NULL} při chybě
    \end{itemize}

    \pause

    \item Zavření: \code{int dlclose(void *handle)}

    \pause

    \item Symbol z knihovny: \code{void *dlsym(void *handle, const char *symbol)}
    \begin{itemize}
      \item \code{handle}: držátko z \code{dlopen}
      \item \code{symbol}: název symbolu (funkce, proměnné, \dots)
      \item Návratová hodnota: ukazatel na začátek symbolu
    \end{itemize}

    \pause

    \item Řetězec chyby: \code{char *dlerror(void)}
  \end{itemize}

\end{frame}

\begin{frame}[fragile]{Načtení dynamické knihovny -- příklad}

  \begin{lstlisting}[aboveskip=2pt,belowskip=2pt]
#include <dlfcn.h>
#include <err.h>
#include <stdio.h>

int main(void) {
  void *h = dlopen("libm.so.6", RTLD_LAZY);
  if (!h)
    errx(1, "dlopen: %s", dlerror());
  \end{lstlisting}
  \begin{uncoverenv}<2->
  \begin{lstlisting}
  double (*kos)(double) = dlsym(h, "cos");
  if (!kos)
    errx(1, "dlsym(cos): %s", dlerror());
  \end{lstlisting}
  \end{uncoverenv}
  \begin{uncoverenv}<3->
  \begin{lstlisting}
  printf("%f\n", kos(0));
  \end{lstlisting}
  \end{uncoverenv}
  \begin{lstlisting}
  dlclose(h);

  return 0;
}
  \end{lstlisting}

\end{frame}

\begin{frame}{Úkol}

  \heading{Načtení dynamické knihovny}
  \begin{enumerate}
    \item Načtěte dynamickou knihovnu
    \begin{itemize}
      \item Zadaná na příkazové řádce
    \end{itemize}

    \item Načtěte symbol
    \begin{itemize}
      \item Zadaný také na příkazové řádce
    \end{itemize}

    \item Zavolejte symbol
    \begin{itemize}
      \item S řetězcem jako parametr
    \end{itemize}
    \item Vyzkoušejte
    \begin{itemize}
      \item Např.\ \cmd{./main libc.so.6 puts}
      \item A \code{printf}
      \item A \code{fclose}?
    \end{itemize}
  \end{enumerate}

\end{frame}

\mysection{Vytváření dynamických knihoven}

\begin{frame}{Pojmenování dynamických knihoven}

  \begin{itemize}
    \item Jméno knihovny (,,soname'')
    \begin{itemize}
      \item \file{libXXX.so.maj_verze}
      \item Např.\ \file{libelf.so.0}
    \end{itemize}

    \item Reálné jméno (soubor)
    \begin{itemize}
      \item \file{libXXX.so.maj_verze.min_verze.release}
      \item Např.\ \file{libelf.so.0.8.13}
    \end{itemize}

    \item Jméno pro linker
    \begin{itemize}
      \item \file{libXXX.so}
      \item Např.\ \file{libelf.so}
    \end{itemize}
  \end{itemize}

\end{frame}

\begin{frame}[fragile]{Vytváření dynamických knihoven}

  \begin{itemize}
    \item Všechen kód musí být nezávislý na pozici, kde bude nahrán
    \begin{itemize}
      \item $\Rightarrow$ nutno překládat i linkovat s \cmd{\-fPIC}, popř.\ 
        lépe s \cmd{\-fpic}
    \end{itemize}

    \pause

    \item Linker by měl znát soname
    \begin{itemize}
      \item $\Rightarrow$ nutno linkovat s \cmd{\-Wl,-soname,jmeno}
    \end{itemize}

    \pause

    \item Linkovat sdíleně
    \begin{itemize}
      \item $\Rightarrow$ nutno linkovat s \cmd{\-shared}
    \end{itemize}

  \end{itemize}

  \begin{block}{Příklad vytvoření knihovny}
  \begin{lstlisting}[language=bash]
gcc -fpic -g -c -Wall a.c
gcc -fpic -g -c -Wall b.c
gcc -fpic -shared -Wl,-soname,libmylib.so.1 -o libmylib.so.1.2.3 a.o b.o -lc
  \end{lstlisting}
  \end{block}

\end{frame}

\begin{frame}{Úkol}

  \heading{Vytvoření dynamické knihovny}
  \begin{enumerate}
    \iffimuni {
      \item Z \file{lib.o} vytvořte dynamickou knihovnu
      \begin{itemize}
	\item \cmd{gcc -shared -Wl,-soname,...}
      \end{itemize}
      \item Knihovnu přilinkujte k \file{main.o} namísto statické
      \begin{itemize}
	\item \cmd{gcc -L. -lX}
      \end{itemize}

      \item Ověte pomocí \cmd{ldd}
      \item Spusťte výsledek
    } \else {
      \item Pracujte v adresáři \file{13} repozitáře
      \item Používejte k překladu/linkování jen příkazovou řádku

      \item Vytvořte dynamickou knihovnu
      \begin{itemize}
	\item Ze zdrojů \file{mylib.*}
	\item Dopište nějakou funkci, která něco vypíše
      \end{itemize}

      \item Ve funkci \code{main} z \file{main.c} ji zavolejte

      \item Slinkujte a vyzkoušejte
      \begin{itemize}
	\item Pomocí \cmd{ldd} ověřte, že se vše povedlo
      \end{itemize}
    } \fi
  \end{enumerate}

\end{frame}
